\section{Introduction}
\label{sec:introduction}


Congruence closure is a important algorithm in computer science. It is used by
(1)~satisfiability modulo theories (SMT) solvers to determine the truth of a
logical formulas~\cite{smt}; and (2)~equality saturation engines to determine
the equivalence of two programs~\cite{eqsat}. Congruence closure determines
whether a set of equalities and disequalities between terms is satisfiable. For
example, given the equalities $a = b$ and $b = c$ and the disequality $f(a) \neq
f(c)$, congruence closure will conclude that this is unsatisfiable since $a = c$
by \emph{transitivity} and thus the two ``parent terms'' are equal ($f(a) =
f(c)$) by \emph{congruence}. To decide this problem in practice, one must
maintain equalities using a union-find data structure and propagate congruences
to parent terms.


% Satisfiability modulo theories (SMT) solvers are used to decided the
% satisfiability of first-order logic formulas. State-of-the-art solvers such as
% Z3~\cite{z3} and cvc5~\cite{cvc5} can reason about a wide variety of
% ``theories'' such as uninterpreted functions, linear arithmetic, bit-vectors,
% and many others.

Both equality saturation engines like egg~\cite{egg} and egglog~\cite{egglog} and SMT solvers like
Z3~\cite{z3} and cvc5~\cite{cvc5} are widely used in a variety of applications,
including software verification, hardware design, and mathematical theorem
proving. Hence, their performance is critical, and much work has gone into
improving their efficiency. One promising avenue for improving the performance
is \emph{parallelization}: utilizing multiple processors to
solve the problem faster. 

To this end, we present a novel parallel congruence closure algorithm. Our
algorithm uses a concurrent union find data structure and a bulk-synchronous
parallel design to efficiently compute the congruence closure of a set of
equalities and disequalities.


Sequential congruence closure algorithms are based on the seminal work by Nelson
and Oppen~\cite{nelson-oppen-cc}, which maintains a worklist of pairs of terms
that are known to be equal. The algorithm repeatedly processes pairs from the
worklist, merging their equivalence classes in the union-find data structure and
adding new pairs to the worklist until no more pairs can be added.

At a high level, our algorithm replaces the classical sequential worklist of
Nelson and Oppen~\cite{nelson-oppen-cc} with a \emph{bulk-synchronous} closure
loop.
Initial equalities are unioned into a lock-free concurrent
union-find~\cite{jayanti16}, marking each affected e-class as \emph{changed}.
The algorithm then proceeds in rounds. In each round we (i)~collect the current
frontier of changed classes in parallel, (ii)~gather every parent e-node whose
canonical signature could have shifted, (iii)~hash each parent's signature under
the current union-find roots, (iv)~group nodes by signature using a parallel
semisort~\cite{semisort}, and (v)~union together every group that witnesses a
new congruence. Each step is a pure data-parallel primitive—\texttt{filter},
\texttt{flat\_map}, \texttt{map}, semisort, and a parallel union pass—built on
\textsc{ParlayLib}'s scheduler and sequence types~\cite{parlaylib}. 

The algorithm terminates when no round discovers a new merge, exactly mirroring the
fixed point of the sequential procedure. 

While congruence closure is P-complete~\cite{ccpcomplete}, we show that our
algorithm can achieve linear speedups on random benchmarks, a synthetic
benchmark suite designed to stress test our algorithm, and a set of equality
saturation benchmarks~\cite{eqsat} from the Herbie project~\cite{herbie}. Our
algorithm efficiently utilizes multiple cores and we discuss what would
be required to integrate it into an existing SMT solver or equality saturation
engines. 